% from ega2-1-fr.tex
\oldpage[II]{74}

$r''=r_{\mathcal{L},\varphi}$, puis $\mathcal{L}_{r''}$ et
$\varphi_{r''}$~; montrons qu'il y a un isomorphisme canonique
$\tau:\mathcal{L}_{r''}\xrightarrow{\sim}\mathcal{L}$ tel que
$\varphi=\tau\circ\varphi_{r''}$~; pour le définir on peut se placer dans le
cas où $Y=\operatorname{Spec}(A)$ et $X=\operatorname{Spec}(B)$ sont affines,
et (avec les notations de (\hyperref[II.4.2.1-fr]{4.2.1}) et de
(\hyperref[II.3.7.2-fr]{3.7.2})) faire correspondre à tout élément
$(x/1)\otimes1$ de $L_{r''}$ (avec $x\in E$) l'élément $v_1(x)c$ de $L$. On
vérifie aussitôt que $\tau$ ne dépend pas du générateur choisi $c$ de $L$~;
comme $v_1$ est par hypothèse surjectif, pour prouver que $\tau$ est un
isomorphisme, il suffit de voir que si $x/1=0$ dans $(S(1))_{(f)}$, on a
$v_1(x)/1=0$ dans $B_g$~; mais la première relation signifie que $f^nx=0$
dans $\mathbf{S}_{n+1}(E)$ pour un certain $n$, et on en déduit que
$v_{n+1}(f^nx)=g^nv_1(x)=0$ dans $B$, d'où la conclusion. Enfin, il est
immédiat que pour deux couples équivalents $(\mathcal{L},\varphi)$,
$(\mathcal{L}',\varphi')$, on a
$r_{\mathcal{L},\varphi}=r_{\mathcal{L}',\varphi'}$.
\end{proof}

En particulier, pour $X=Y$~:

\begin{theorem}[4.2.4]
\label{II.4.2.4-fr}
L'ensemble des $Y$-sections de $\mathbf{P}(\mathcal{E})$ est en
correspondance biunivoque canonique avec l'ensemble des
sous-$\mathcal{O}_Y$-Modules quasi-cohérents $\mathcal{F}$ de $\mathcal{E}$
tels que $\mathcal{E}/\mathcal{F}$ soit inversible.
\end{theorem}

On notera que cette propriété correspond à la définition classique de
l'\og espace projectif\fg{} comme ensemble des hyperplans d'un espace
vectoriel (le cas classique correspondant à $Y=\operatorname{Spec}(K)$, où
$K$ est un corps, et $\mathcal{E}=\widetilde{E}$, $E$ étant un $K$-espace
vectoriel de dimension finie~; les $\mathcal{F}$ ayant la propriété énoncée
dans (\hyperref[II.4.2.4-fr]{4.2.4}) correspondent alors aux hyperplans de
$E$, et on sait d'autre part que les $Y$-sections de
$\mathbf{P}(\mathcal{E})$ sont alors les points de
$\mathbf{P}(\mathcal{E})$ rationnels sur $K$
(\hyperref[I.3.4.5-fr]{I, 3.4.5})).

\begin{remark}[4.2.5]
\label{II.4.2.5-fr}
Comme il y a correspondance biunivoque canonique entre $Y$-morphismes de $X$
dans $P$ et leurs morphismes graphes, $X$-sections de $P\times_YX$
(\hyperref[I.3.3.14-fr]{I, 3.3.14}), on voit qu'inversement
(\hyperref[II.4.2.3-fr]{4.2.3}) peut se déduire de
(\hyperref[II.4.2.4-fr]{4.2.4}). Désignons par
$\operatorname{Hyp}_Y(X,\mathcal{E})$ l'ensemble des
sous-$\mathcal{O}_X$-Modules quasi-cohérents $\mathcal{F}$ de
$\mathcal{E}\otimes_Y\mathcal{O}_X=q^*(\mathcal{E})$ tels que
$q^*(\mathcal{E})/\mathcal{F}$ soit un $\mathcal{O}_X$-Module inversible. Si
$g:X'\to X$ est un $Y$-morphisme, il résulte du fait que $g^*$ est exact à
droite que l'on a
\[
  g^*(q^*(\mathcal{E})/\mathcal{F})=
  g^*(q^*(\mathcal{E}))/g^*(\mathcal{F}),
\]
donc que ce dernier faisceau est inversible, et par suite
$\operatorname{Hyp}_Y(X,\mathcal{E})$ est un foncteur contravariant dans la
catégorie des $Y$-préschémas. On peut alors interpréter le th.
(\hyperref[II.4.2.4-fr]{4.2.4}) comme définissant un isomorphisme canonique
des foncteurs $\operatorname{Hom}_Y(X,\mathbf{P}(\mathcal{E}))$ et
$\operatorname{Hyp}_Y(X,\mathcal{E})$ contravariants en $X$ variable dans la
catégorie des $Y$-préschémas. Ceci fournit aussi une caractérisation du fibré
projectif $P=\mathbf{P}(\mathcal{E})$ par la propriété universelle suivante,
plus proche de l'intuition géométrique que les constructions des \S\S~2 et
3~: pour tout morphisme $q:X\to Y$ et tout $\mathcal{O}_X$-Module inversible
$\mathcal{L}$, quotient de $\mathcal{E}\otimes_{\mathcal{O}_Y}\mathcal{O}_X$,
il existe un $Y$-morphisme unique $r:X\to P$ tel que
$\mathcal{L}=r^*(\mathcal{O}_P(1))$.

On verra plus tard comment on peut de la même manière définir, entre autres,
les schémas \og grassmanniennes\fg{}.
\end{remark}

\begin{corollary}[4.2.6]
\label{II.4.2.6-fr}
Supposons que tout $\mathcal{O}_Y$-Module inversible soit trivial
(\hyperref[I.2.4.8-fr]{I, 2.4.8}). Soit $V$ le groupe
$\operatorname{Hom}_{\mathcal{O}_Y}(\mathcal{E},\mathcal{O}_Y)$, considéré
comme module sur l'anneau $A=\Gamma(Y,\mathcal{O}_Y)$, et soit $V^*$ le
sous-ensemble de $V$ formé des homomorphismes surjectifs. Alors l'ensemble
des $Y$-sections de $\mathbf{P}(\mathcal{E})$ s'identifie canoniquement à
$V^*/A^*$, où $A^*$ est le groupe des unités de $A$.
\end{corollary}

